\documentclass[10pt, letterpaper]{article}


% Packages:
\usepackage[
    ignoreheadfoot, % set margins without considering header and footer
    top=2 cm, % seperation between body and page edge from the top
    bottom=2 cm, % seperation between body and page edge from the bottom
    left=2 cm, % seperation between body and page edge from the left
    right=2 cm, % seperation between body and page edge from the right
    footskip=1.0 cm, % seperation between body and footer
    % showframe % for debugging 
]{geometry} % for adjusting page geometry
\usepackage[explicit]{titlesec} % for customizing section titles
\usepackage{tabularx} % for making tables with fixed width columns
\usepackage{array} % tabularx requires this
\usepackage[dvipsnames]{xcolor} % for coloring text
\definecolor{primaryColor}{RGB}{0, 79, 144} % define primary color
\usepackage{enumitem} % for customizing lists
\usepackage{fontawesome5} % for using icons
\usepackage{amsmath} % for math
\usepackage[
    pdftitle={简历-李梦成},
    pdfauthor={李梦成},
    pdfcreator={LaTeX with RenderCV},
    colorlinks=true,
    urlcolor=primaryColor
]{hyperref} % for links, metadata and bookmarks
\usepackage[pscoord]{eso-pic} % for floating text on the page
\usepackage{calc} % for calculating lengths
\usepackage{bookmark} % for bookmarks
\usepackage{lastpage} % for getting the total number of pages
\usepackage{changepage} % for one column entries (adjustwidth environment)
\usepackage{paracol} % for two and three column entries
\usepackage{ifthen} % for conditional statements
\usepackage{needspace} % for avoiding page brake right after the section title
\usepackage{iftex} % check if engine is pdflatex, xetex or luatex

\usepackage{ctex}
\setCJKmainfont{SimSun}[
AutoFakeBold = 3,
ItalicFont   = KaiTi,
]%
\setCJKsansfont{SimHei}[AutoFakeBold = 3]%
\setCJKmonofont{FangSong}%
\setCJKfamilyfont{zhsong}{SimSun}[AutoFakeBold = 3]%
\setCJKfamilyfont{zhhei}{SimHei}[AutoFakeBold = 3]%
\setCJKfamilyfont{zhkai}{KaiTi}%
\setCJKfamilyfont{zhfs}{FangSong}%

\usepackage{fontspec}
\setmainfont{Times New Roman}          % 先设置主字体
% \usepackage[default, type1]{sourcesanspro} % 后加载无衬线字体

% Some settings:
\AtBeginEnvironment{adjustwidth}{\partopsep0pt} % remove space before adjustwidth environment
\pagestyle{empty} % no header or footer
\setcounter{secnumdepth}{0} % no section numbering
\setlength{\parindent}{0pt} % no indentation
\setlength{\topskip}{0pt} % no top skip
\setlength{\columnsep}{0.15cm} % set column seperation
\makeatletter
\let\ps@customFooterStyle\ps@plain % Copy the plain style to customFooterStyle
\patchcmd{\ps@customFooterStyle}{\thepage}{
    \color{gray}\textit{\small 李梦成 - Page \thepage{} of \pageref*{LastPage}}
}{}{} % replace number by desired string
\makeatother
\pagestyle{customFooterStyle}

\titleformat{\section}{
    % make the font size of the section title large and color it with the primary color
    \Large\color{primaryColor}
}{
}{
}{
    % print bold title, give 0.15 cm space and draw a line of 0.8 pt thickness
    % from the end of the title to the end of the body
    \textbf{#1}\hspace{0.15cm}\titlerule[0.8pt]\hspace{-0.1cm}
}[] % section title formatting

\titlespacing{\section}{
    % left space:
    -1pt
}{
    % top space:
    0.3 cm
}{
    % bottom space:
    0.2 cm
} % section title spacing

% \renewcommand\labelitemi{$\vcenter{\hbox{\small$\bullet$}}$} % custom bullet points
\newenvironment{highlights}{
    \begin{itemize}[
        topsep=0.10 cm,
        parsep=0.10 cm,
        partopsep=0pt,
        itemsep=0pt,
        leftmargin=0.4 cm + 10pt
    ]
}{
    \end{itemize}
} % new environment for highlights

\newenvironment{highlightsforbulletentries}{
    \begin{itemize}[
        topsep=0.10 cm,
        parsep=0.10 cm,
        partopsep=0pt,
        itemsep=0pt,
        leftmargin=10pt
    ]
}{
    \end{itemize}
} % new environment for highlights for bullet entries


\newenvironment{onecolentry}{
    \begin{adjustwidth}{
        0.2 cm + 0.00001 cm
    }{
        0.2 cm + 0.00001 cm
    }
}{
    \end{adjustwidth}
} % new environment for one column entries

\newenvironment{twocolentry}[2][]{
    \onecolentry
    \def\secondColumn{#2}
    \setcolumnwidth{\fill, 4.8 cm}
    \begin{paracol}{2}
}{
    \switchcolumn \raggedleft \secondColumn
    \end{paracol}
    \endonecolentry
} % new environment for two column entries

\newenvironment{threecolentry}[3][]{
    \onecolentry
    \def\thirdColumn{#3}
    \setcolumnwidth{1 cm, \fill, 4.5 cm}
    \begin{paracol}{3}
    {\raggedright #2} \switchcolumn
}{
    \switchcolumn \raggedleft \thirdColumn
    \end{paracol}
    \endonecolentry
} % new environment for three column entries

\newenvironment{header}{
    \setlength{\topsep}{0pt}\par\kern\topsep\centering\color{primaryColor}\linespread{1.5}
}{
    \par\kern\topsep
} % new environment for the header

% \newcommand{\dateRange}[2]{
%     {#1 -- \makebox[\widthof{#1}][s]{#2}}
% }
% \newcommand{\dateRange}[2]{
%     {\makebox[\widthof{2018年12月}][s]{#1} -- \makebox[\widthof{2018年12月}][s]{#2}}
% }
\newcommand{\dateRange}[2]{
    {\makebox[\widthof{2018年12月}][c]{#1} -- \makebox[\widthof{2018年12月}][c]{#2}}
}

\newcommand{\placelastupdatedtext}{% \placetextbox{<horizontal pos>}{<vertical pos>}{<stuff>}
  \AddToShipoutPictureFG*{% Add <stuff> to current page foreground
    \put(
        \LenToUnit{\paperwidth-2 cm-0.2 cm+0.05cm},
        \LenToUnit{\paperheight-1.0 cm}
    ){\vtop{{\null}\makebox[0pt][c]{
        \small\color{gray}\textit{Last updated in September 2024}\hspace{\widthof{Last updated in September 2024}}
    }}}%
  }%
}%

% save the original href command in a new command:
\let\hrefWithoutArrow\href

% new command for external links:
\renewcommand{\href}[2]{\hrefWithoutArrow{#1}{\ifthenelse{\equal{#2}{}}{ }{#2 }\raisebox{.15ex}{\footnotesize \faExternalLink*}}}


\begin{document}
    \newcommand{\AND}{\unskip
        \cleaders\copy\ANDbox\hskip\wd\ANDbox
        \ignorespaces
    }
    \newsavebox\ANDbox
    \sbox\ANDbox{}

    % \placelastupdatedtext
    \begin{header}
        \fontsize{30 pt}{30 pt}
        \textbf{李梦成}
 
        \vspace{0.3 cm}

        \normalsize
        \mbox{{\footnotesize\faMapMarker*}\hspace*{0.13cm}北京}%
        \kern 0.25 cm%
        \AND%
        \kern 0.25 cm%
        \mbox{\hrefWithoutArrow{mailto:mengcheng_li@163.com}{{\footnotesize\faEnvelope[regular]}\hspace*{0.13cm}mengcheng\_li@163.com}}%
        \kern 0.25 cm%
        \AND%
        \kern 0.25 cm%
        \mbox{\hrefWithoutArrow{tel:18401653211}{{\footnotesize\faPhone*}\hspace*{0.13cm}18401653211}}%
        \kern 0.25 cm%
        \AND%
        \kern 0.25 cm%
        \mbox{\hrefWithoutArrow{https://dw1010.github.io/}{{\footnotesize\faLink}\hspace*{0.13cm}个人主页}}%
        \kern 0.25 cm%
        \AND%
        \kern 0.25 cm%
        % \mbox{\hrefWithoutArrow{https://linkedin.com/in/yourusername}{{\footnotesize\faLinkedinIn}\hspace*{0.13cm}yourusername}}%
        % \kern 0.25 cm%
        % \AND%
        % \kern 0.25 cm%
        \mbox{\hrefWithoutArrow{https://github.com/Dw1010}{{\footnotesize\faGithub}\hspace*{0.13cm}github}}%
    \end{header}

    \vspace{0.3 cm - 0.3 cm}

    % \section{Welcome to RenderCV!}

    %     \begin{onecolentry}
    %         \href{https://rendercv.com}{RenderCV} is a LaTeX-based CV/resume version-control and maintenance app. It allows you to create a high-quality CV or resume as a PDF file from a YAML file, with \textbf{Markdown syntax support} and \textbf{complete control over the LaTeX code}.
    %     \end{onecolentry}

    %     \vspace{0.2 cm}

    %     \begin{onecolentry}
    %         The boilerplate content was inspired by \href{https://github.com/dnl-blkv/mcdowell-cv}{Gayle McDowell}.
    %     \end{onecolentry}


    
    % \section{Quick Guide}

    % \begin{onecolentry}
    %     \begin{highlightsforbulletentries}


    %     \item Each section title is arbitrary and each section contains a list of entries.

    %     \item There are 7 unique entry types: \textit{BulletEntry}, \textit{TextEntry}, \textit{EducationEntry}, \textit{ExperienceEntry}, \textit{NormalEntry}, \textit{PublicationEntry}, and \textit{OneLineEntry}.

    %     \item Select a section title, pick an entry type, and start writing your section!

    %     \item \href{https://docs.rendercv.com/user_guide/}{Here}, you can find a comprehensive user guide for RenderCV.


    %     \end{highlightsforbulletentries}
    % \end{onecolentry}

    % \section{\textbf{研究方向}}

    % \begin{onecolentry}
    %     \begin{highlightsforbulletentries}
        
    %     \item 人手的运动捕捉与三维重建

    %     \item 人手交互与运动生成

    %     \item 人手化身生成与重光照
        
    %     \item 生成式图像复原与超分 

    %     \end{highlightsforbulletentries}
    % \end{onecolentry}

    \section{\textbf{教育经历}}
        \begin{twocolentry}
            {
            \dateRange{2014年9月}{2018年6月}
            GPA 89/100
            }
            {
            \textbf{清华大学物理系}，数学物理基础科学

            本科
            }
        \end{twocolentry}

        \begin{twocolentry}
            {
            \dateRange{2018年9月}{2025年6月}
            GPA 3.73/4
            }
            {
            \textbf{清华大学自动化系}，控制理论与工程

            博士（导师：\hrefWithoutArrow{https://www.au.tsinghua.edu.cn/info/1080/3160.htm}{刘烨斌}）

            \textbf{主要研究方向：}
            \begin{highlightsforbulletentries}
                \item 人手的运动捕捉与三维重建

                \item 人手交互与运动生成

                \item 人手化身生成与重光照
            \end{highlightsforbulletentries}

            }
        \end{twocolentry}
        
        % \begin{threecolentry}{\textbf{BS}}{
        %     Sept 2000 – May 2005
        % }
        %     \textbf{University of Pennsylvania}, Computer Science
        %     \begin{highlights}
        %         \item GPA: 3.9/4.0 (\href{https://example.com}{a link to somewhere})
        %         \item \textbf{Coursework:} Computer Architecture, Comparison of Learning Algorithms, Computational Theory
        %     \end{highlights}
        % \end{threecolentry}

    \section{\textbf{实习经历}}
        \begin{twocolentry}
            {
            \dateRange{2018年12月}{2019年3月}
            中国北京
            }
            {
            \textbf{北京市商汤科技开发有限公司}
            
            实习

            \begin{highlightsforbulletentries}
                \item 工作内容：主要负责数字人驱动，包括参数人体模板（SMPL）的驱动以及对人体unity数字人的运动迁移
            \end{highlightsforbulletentries}
            }
        \end{twocolentry}

        \begin{twocolentry}
            {
            \dateRange{2020年7月}{2020年9月}
            中国北京
            }
            {
            \textbf{北京国成万通信息科技公司}

            实习

            \begin{highlightsforbulletentries}
                \item 工作内容：主要负责人手抓取物体的物理仿真以及对机械臂抓取方式的迁移
            \end{highlightsforbulletentries}
            }
        \end{twocolentry}
    
    \section{\textbf{工作经历}}
        \begin{twocolentry}
            {
            \dateRange{2025年6月}{至今}
            中国北京
            }
            {
            \textbf{北京荣耀终端有限公司}
            
            影像算法部高级算法工程师

            \begin{highlightsforbulletentries}
                \item 工作内容：图片复原与超分，尤其是基于生成式算法的长焦小文字图片的复原与超分
            \end{highlightsforbulletentries}
            }
        \end{twocolentry}

    \section{\textbf{项目经历}}
        \begin{onecolentry}
            \begin{highlightsforbulletentries}
                \item \textbf{[主要负责人] 基于单RGB图片的双手重建}\par
                实现了一个实时的基于单张RGB图片的紧密双手交互重建系统。
                主要创新点为图卷积网络和注意力机制的结合，解决了双手紧密交互下的复杂遮挡问题。
                相关论文发表在CVPR2022会议上，并在会议上进行了口头汇报（oral）。论文目前已经获得130+引用量。\par
                项目主页：\hrefWithoutArrow{https://dw1010.github.io/project/IntagHand/Intaghand.html}{https://dw1010.github.io/project/IntagHand/Intaghand.html}
                \item \textbf{[主要负责人] 基于图扩散模型的多模态人手生成}\par
                实现了一个多模态的基于图卷积和扩散模型的多模态人手生成网络。
                能够在不同模态的条件下实现人手三维模型的生成。
                相关论文发表在CVPR2024会议上，并被选为亮点论文（HighLight）。\par
                项目主页：\hrefWithoutArrow{https://dw1010.github.io/project/HHMR/HHMR.html}{https://dw1010.github.io/project/HHMR/HHMR.html}
                \item \textbf{[主要负责人] 基于三维高斯泼溅的可重光照人手化身}\par
                实现了一个三维高斯泼溅的可重光照人手化身生成算法。
                该算法将基于物理的渲染方法引入三维高斯泼溅，实现了环境光照和人手本征纹理的分解，从而实现了高精度的可重光照人手化身。\par
                Demo展示：\hrefWithoutArrow{https://dw1010.github.io/}{https://dw1010.github.io/}
                \item \textbf{[参与] 轻量级多人全身动作捕捉}\par
                实现了一个轻量化的无标记多人体实时的全身动作捕捉系统。
                能够从多视角的RGB视频同时捕捉多人的身体、人手、脸部的动作。
                相关论文发表在ICCV2021会议上。\par
                项目主页：\hrefWithoutArrow{https://www.liuyebin.com/lwtotalcap/lwtotalcap.html}{https://www.liuyebin.com/lwtotalcap/lwtotalcap.html}
                \item \textbf{[参与] 基于单RGB图片的高精度全身姿态回归}\par
                实现了一个高精度的全身姿态回归算法PyMaf-X。
                其通过金字塔网格对齐反馈循环提升单目图像全身模型恢复的准确性和自然度。
                相关论文发表在TPAMI2023上。\par
                项目主页：\hrefWithoutArrow{https://www.liuyebin.com/pymaf-x/}{https://www.liuyebin.com/pymaf-x/}
                \item \textbf{[参与、指导] 基于单RGB手物交互图片的人手和物体重建}\par
                实现了从人手抓握物体的单张RGB图片中同时重建人手和物体的算法。
                其中物体的重建利用了人手的稠密顶点接触信息来估计物体的SDF场。
                相关论文发表在AAAI2024上。\par
                项目主页：\hrefWithoutArrow{https://junxinghu.github.io/projects/hoi.html}{https://junxinghu.github.io/projects/hoi.html}
                \item \textbf{[参与] 人手与铰链物体交互的运动序列生成}\par
                实现了一种人手与铰链物体运动交互序列生成的算法。
                基于扩散生成模型，首先生成稠密的手物交互区域信息，然后通过交互约束的残差迭代生成交互过程中的手部姿态，实现了对刚体和铰接物体的单/双手逼真操控。
                相关论文发表在TPAMI2025上。\par
                项目主页：\hrefWithoutArrow{https://jiajunzhang16.github.io/manidext/}{https://jiajunzhang16.github.io/manidext/}
                \item \textbf{[参与、指导] 基于单RGB视频序列的双手重建}\par
                实现了一个基于RGB视频序列的双手重建大模型。
                分别输入两只手的局部裁切图片，并通过基于二维相对位置编码实现相对三维空间位置的估计，提高了双手交互时相互深度估计的稳定性，
                在重建紧密交互双手上实现了超越以往方法的稳定性和准确度。相关论文发表在TOG2026上。\par
                项目主页：\hrefWithoutArrow{https://omnihand.github.io/}{https://omnihand.github.io/}
                \item \textbf{[主要负责人] 长焦小文字图片超分}\par
                实现了一个生成式的小文字超分算法。
                构建了一套文字数据渲染管线，使用文生图模型生成图片+OCR与vlm模型检验筛选数据。
                使用clip模型对齐标准文字图片和文字笔画结构的特征分布，得到语义丰富的文字token。
                此外，对于细小文字笔画设计了一套流模型生成算法，能够在保持文字结构的同时提升小文字的清晰度和可读性。\par
            \end{highlightsforbulletentries}
        \end{onecolentry}

    
    \section{\textbf{学术论文}}

        \begin{onecolentry}
            \begin{highlightsforbulletentries}
                \item \textbf{Mengcheng Li}, Liang An, Tao Yu, Yangang Wang, Feng Chen, Yebin Liu.\par
                \textit{Neural hand reconstruction using an RGB image.}\par
                Virtual Reality \& Intelligent Hardware, 2020, 2(3): 276-289 DOI:10.1016/j.vrih.2020.05.001
                \item \textbf{Mengcheng Li}, Liang An, Hongwen Zhang, Lianpeng Wu,, Feng Chen, Tao Yu, Yebin Liu.\par
                \textit{Interacting Attention Graph for Single Image Two-Hand Reconstruction.}\par
                IEEE/CVF Conference on Computer Vision and Pattern Recognition (CVPR), 2022 (\textbf{Oral Paper}).
                \item \textbf{Mengcheng Li}, Hongwen Zhang, Yuxiang Zhang, Ruizhi Shao, Tao Yu, Yebin Liu.\par
                \textit{HHMR: Holistic Hand Mesh Recovery by Enhancing the Multimodal Controllability of Graph Diffusion Models.}\par
                IEEE/CVF Conference on Computer Vision and Pattern Recognition (CVPR), 2024 (\textbf{HighLight Paper}).
                \item Yuxiang Zhang, Zhe Li, Tao Yu,\textbf{Mengcheng Li}, Liang An, Yebin Liu.\par
                \textit{Light-weight Multi-person Total Capture Using Sparse Multi-view Cameras.}\par
                IEEE/CVF International Conference on Computer Vision (ICCV), 2021.
                \item Hongwen Zhang, Yating Tian, Yuxiang Zhang, \textbf{Mengcheng Li}, Liang An, Zhenan Sun, Yebin Liu.\par
                \textit{PyMAF-X: Towards Well-aligned Full-body Model Regression from Monocular Images.}\par 
                IEEE Transactions on Pattern Analysis and Machine Intelligence (TPAMI), 2023.
                \item Junxing Hu, Hongwen Zhang, Zerui Chen, \textbf{Mengcheng Li}, Yunlong Wang, Yebin Liu, Zhenan Sun.\par
                \textit{Learning Explicit Contact for Implicit Reconstruction of Hand-held Objects from Monocular Images.}\par
                AAAI Conference on Artificial Intelligence (AAAI), 2024.
                \item Jiajun Zhang, Yuxiang Zhang, Liang An, \textbf{Mengcheng Li}, Hongwen Zhang, Zonghai Hu, Yebin Liu.\par
                \textit{ManiDext: Hand-Object Manipulation Synthesis via Continuous Correspondence Embeddings and Residual-Guided Diffusion.}\par
                IEEE Transactions on Pattern Analysis and Machine Intelligence (TPAMI), 2025
                \item Dixuan Lin, Yuxiang Zhang, \textbf{Mengcheng Li}, Yebin Liu, Wei Jing, Qi Yan, Qianying Wang, Hongwen Zhang.\par
                \textit{4DHands: Reconstructing Interactive Hands in 4D with Transformers.}\par
                ACM Transactions on Graphics (TOG), 2026.
            \end{highlightsforbulletentries}
        \end{onecolentry}


    \section{\textbf{发明专利}}
        \begin{onecolentry}
            \begin{highlightsforbulletentries}
                \item (ZL201910464503.7) \quad 基于深度学习的单图像的人手动作与形状重建方法及装置 （已授权）
                \item (ZL201910646553.7) \quad 手部运动重建方法和装置 （已授权）
                \item (CN118015167A) \quad 基于光流的单视频双手动态三维运动捕捉方法和装置 （已公开）
            \end{highlightsforbulletentries}
        \end{onecolentry}

    \section{\textbf{个人技能}}
        \begin{onecolentry}
            \textbf{语言:} 汉语，英语
        \end{onecolentry}
        \begin{onecolentry}
            \textbf{编程语言:} C，C++，Python，CUDA
        \end{onecolentry}
        \begin{onecolentry}
            \textbf{开发工具:} PyTorch，Pytorch3D，OpenGL，Nvdiffrast
        \end{onecolentry}
\end{document}